% Generated draft: requires mathematical and source review.
\documentclass[11pt]{article}
\usepackage[margin=25mm]{geometry}
\usepackage{fontspec}
\setmainfont{Latin Modern Roman}
\usepackage{amsmath,amssymb,mathtools,unicode-math}
\setmathfont{Latin Modern Math}
\usepackage{enumitem,needspace,xurl,hyperref}
\setlength{\parindent}{0pt}
\setlength{\parskip}{5pt}
\setlength{\emergencystretch}{3em}
\newcommand{\N}{\mathbb{N}}
\newcommand{\Z}{\mathbb{Z}}
\newcommand{\Q}{\mathbb{Q}}
\newcommand{\R}{\mathbb{R}}
\newcommand{\C}{\mathbb{C}}
\newcommand{\F}{\mathbb{F}}
\newcommand{\A}{\mathbb{A}}
\newcommand{\PP}{\mathbb{P}}
\newcommand{\E}{\mathbb{E}}
\newcommand{\eps}{\varepsilon}
\newcommand{\dd}{\,\mathrm{d}}
\begin{document}
\section*{UnsolvedMath ID: 20000467}
\textbf{Source problem code:} catalogue entry
\subsection{Definitions and mathematical statement}
For \(P=[x_0:\cdots:x_m]\in\PP^m(\Q)\), use primitive integral coordinates and the height \(H(P)=\max_i|x_i|\). Let \(C\subset\PP^2_{\Q}\) be the plane curve defined by a homogeneous degree-\(d\) polynomial, and put \(N_C(B)=\#\{P\in C(\Q):H(P)\le B\}\). The supplied interpretation of Bogomolov's damaged question asks to replace Heath-Brown's method by an Arakelov method for a curve bound of the form
\[
 N_C(B)\ll B^{\alpha+\epsilon},\qquad \alpha\le\frac2d,
\]
and to seek analogues for surfaces and higher dimensions. Here \(\epsilon>0\); the dependence of the implicit constant and the needed irreducibility/smoothness hypotheses are not stated in the extracted source. For example, a degree-\(d\) equation with a rational line component cannot be treated as an arbitrary integral degree-\(d\) curve for this bound. Thus the standard geometrically integral curve setting must be specified when using the reported theorem, rather than silently imposed on all polynomials in the source.

The supplied report says that Chen's published Arakelov slope-method covering theorem supplies the desired methodological replacement for the degree-\(d\) curve bound. Its separate candidate exponent-transfer lemma is: for every smooth degree-\(d\) surface \(X\subset\PP^3_{\Q}\), every integer \(r\) with \(0\le r\le d-2\), and every \(\epsilon>0\), let \(U_r\) be the locus left after deleting all geometric integral curves on \(X_{\overline\Q}\) of degree at most \(r\). Then
\[
 N_{U_r}(B)=\#\{P\in X(\Q):P\in U_r,\ H(P)\le B\}
 \ll_{d,r,\epsilon} B^{3/\sqrt d+2/(r+1)+\epsilon}.
\]
The description of \(U_r\) is a point-set deletion; the record does not assert that the union being deleted is finite. At \(r=d-2\), the reported exponent is \(3/\sqrt d+2/(d-1)+\epsilon\), identified by the report with the Heath-Brown surface exponent recorded in the surrounding AIM discussion. The report says that this uses Chen's covering theorem and curve corollary, including treatment of intersection components integral over \(\Q\) but geometrically reducible. Their precise statements and proof details are not supplied.

A distinct refinement replacing the \(\epsilon\)-loss by a power of \(\log B\) is explicitly conditional on Chiu's preprint theorem as stated in \texttt{arXiv:2412.05205v1}, combined with the Binyamini--Cluckers--Kato curve bound. No exponent of the logarithm or full hypotheses of that conditional covering theorem are given. This conditional refinement is not merged with the report's asserted unconditional \(\epsilon\)-bound.
\subsection{Short English statement}
Seek an Arakelov-theoretic method for the degree-dependent rational-point curve bound and analogous surface estimates. The supplied report attributes the curve replacement to Chen and gives a surface bound after deleting low-degree geometric curves; a stronger logarithmic refinement is explicitly conditional on a named preprint version.
\subsection{Sources}
\begin{itemize}
\item Local UnsolvedMath registry, permanent ID \(20000467\), source code \texttt{catalogue entry}; AIM algebraic number theory corpus, \texttt{aim-algebraic-number-theory-notes.json}, notes[118], stored source item 3. Attribution: Ulam AI / UnsolvedMath Contributors, CC BY 4.0.
\item AIM \emph{Rational and integral points on higher dimensional varieties}, attributed to Bogomolov: \url{https://aimath.org/WWN/qptsurface2/qptsurface2.pdf}. The supplied verification reports a split question number ``3''/``4'', displaced exponents and inequalities, and navigation text ``Problem/''; it refers to official HTML and PDF page 55 but stops before the announced recovered text. The split number suggests 34, while the stored item remains 3; no metadata is overwritten.
\item The report names Chen's published slope-method covering theorem and curve corollary, Chiu's \texttt{arXiv:2412.05205v1}, and the Binyamini--Cluckers--Kato curve bound. No full citations or theorem statements for these works are supplied; the original literature field is empty. The source's request names Heath-Brown and Arakelov.
\item Registry dated August 14, 2026: statement exact, status \texttt{partially\_solved}, classification \texttt{PARTIAL-PROGRESS}; attempt 1 labelled partial result, rigor proved, novelty confidence low. The exponent-transfer lemma and conditional refinement are machine-generated research aids, not peer reviewed, requiring independent verification before citation.
\end{itemize}
\end{document}
